\cvsection{Professional Experience}
\begin{cventries}
  \cventry
  {Senior Data Scientist / Lead Data Engineer / Senior Technical Manager (Maritime)}
  {SynMax}
  {London Area, United Kingdom}
  {May 2023 - Present}
  {
    \begin{cvitems}
      \item {Spearheaded the development and rollout of a new agentic AI product positioned as the next-generation evolution of SynMax's legacy UI platform, shaping one of the company's most strategically important product transitions.}
      \item {Drove a step change in developer velocity across engineering and data science by embedding AI-enabled ways of working into delivery, accelerating execution and raising the technical ambition of the organisation.}
      \item {Reorganised Python and frontend engineering into a full-stack delivery model, creating clearer ownership, tighter execution and a stronger foundation for cross-functional product delivery at scale.}
      \item {Manage Theia, a maritime domain awareness platform that fuses 30 million square km of satellite imagery with other operational data sources, including 500 million AIS messages per day, to track vessels at scale.}
      \item {Oversee an extensive cloud footprint spanning 200 microservices, with architectural ownership across distributed services, cloud-native data platforms and AI/ML workflows.}
      \item {Built large-scale pipelines for AIS, satellite and sensor data ingestion, transformation, storage and querying, enabling high-volume maritime analytics and operational tracking.}
      \item {Redesigned storage architecture using a hybrid PostgreSQL plus parquet/duckdb approach, cutting annual infrastructure cost from \textasciitilde\$300k to \textasciitilde\$30k while improving analytics speed, backend security and deployment workflows.}
      \item {Built agentic AI systems using LLM-driven and autonomous workflows for text-based intelligence processing and mission-critical maritime applications.}
      \item {Led deep learning and computer vision work on optical satellite imagery for object detection, classification and environmental monitoring, alongside wider geospatial ML and vessel behaviour analytics.}
      \item {Progressed through senior data science, lead data engineering and senior technical management responsibilities while remaining deeply hands-on across architecture, platform engineering, AI delivery and product direction.}
      \item {Lead a multidisciplinary organisation delivering a geospatial AI platform for maritime domain awareness, with responsibility across ingest, data processing, ML systems, frontend and customer-facing data delivery.}
      \item {Manage through 5 tech leads with overall responsibility for 28 staff across data science, data engineering, data annotation, frontend, Python development, AI development and GCP infrastructure.}
      \item {Act as a subject matter expert on AIS data, advising product, engineering and business teams on tracking algorithms, spoofing detection, dark ships, anomalous maritime behaviour and geospatial trends.}
      \item {Led DevOps and infrastructure work across AWS and GCP, including CI/CD, observability, IAM, networking, containerised deployment, Terraform and Kubernetes-based workflows.}
      \item {Support business development and strategic delivery by translating technical capability into customer value propositions, proposals and client-facing technical direction.}
    \end{cvitems}
  }
  \cventry
  {Data Science Team Lead}
  {Fraym}
  {Washington, DC, USA}
  {Oct 2022 - Jun 2023}
  {
    \begin{cvitems}
      \item {Led research data science work focused on stacked ensemble machine learning models for spatial interpolation, forecasting and uncertainty quantification.}
      \item {Managed a small research data science team and worked closely with machine learning engineers and data engineers to translate methodology into product-ready systems.}
      \item {Provided technical leadership on model development for geospatial products, balancing hands-on research with delivery-oriented collaboration.}
    \end{cvitems}
  }
  \cventry
    {Lead Data Scientist / Data Scientist}
    {Geollect}
    {Bristol, UK}
    {Sep 2020 - Oct 2022}
    {
      \begin{cvitems}
        \item {Led a team of data scientists and engineers building SaaS applications, scaling data products and delivering external consultancy across maritime and geospatial intelligence.}
        \item {Built and maintained ETL pipelines, Python services, databases and APIs on AWS using technologies including Terraform, Kinesis, S3, EC2, ECS, DynamoDB and RDS.}
        \item {Served as the company's lead data science consultant for defence, maritime and public-sector customers including work involving AIS analytics, anomaly detection, acoustic data, satellite imagery and intelligence investigations.}
        \item {Developed machine learning solutions for vessel behaviour classification, anomaly detection and AIS-enriched vessel imagery, including combining AIS with satellite imagery for operational workflows.}
        \item {Delivered traditional analytics and visualisation work as well, including digital port assessment products, static reporting and dashboard-style outputs for customer use.}
      \end{cvitems}
    }
  \cventry
    {Postdoctoral Researcher}
    {University of Bath}
    {Bath, UK}
    {Sep 2019 - Sep 2020}
    {
      \begin{cvitems}
        \item {Developed machine learning models to aid in the discovery of new energy materials.}
      \end{cvitems}
    }
  \cventry
    {Postgraduate Teaching Assistant}
    {University of Bath}
    {Bath, UK}
    {Sep 2016 - Jan 2020}
    {
      \begin{cvitems}
        \item {Taught computational chemistry, programming for chemistry, drug discovery, advanced practical chemistry and maths for chemists.}
      \end{cvitems}
    }
\end{cventries}
